\begin{claim}{C.19}{Proved under stated hypotheses}{Coproduct preservation on disk coefficients also closes the raw gate}\label{C19}
Assume the normalized corollary hypotheses of Claim~\ref{C6} and degreewise purity of every $F(\alpha\otimes X)$. Let $F^{(0)}$ be any exact cocontinuous resolving functor computing $L\varphi$, for example $F$ itself. On the category $\mathcal A=\prod_{m\geq1}\Mod k$ define
\[
 DQ=\prod_{m\geq1}D^m(Q_m),\qquad E=F^{(0)}GD.
\]
If $E$ preserves small coproducts, then the raw adjoints $F,G$ induce quasi-inverse equivalences over every commutative ring $k$.

A sufficient hypothesis is that $G$ preserves coproducts. This holds if each $F(h_a)$ is a finitely generated projective object of $\Ch_{\geq0}(k)$, equivalently a finite direct sum, up to retract, of degree-zero spheres and disks on finitely generated projective modules. We make no finiteness assumption on the number of generators $a$, and no inference that this compactness holds for an arbitrary projective resolution.
\end{claim}
\begin{proof}
The functor $D$ is exact and preserves colimits: only two disk contributions occur in any chain degree. Since $F^{(0)}$ and $G$ are exact, so is $E$. Its assumed preservation of coproducts therefore implies preservation of all colimits. Let $e_m:\Mod k\to\mathcal A$ insert a module in coordinate $m$ and zero elsewhere, and put $K_m=E(e_m k)$. Writing a module as the cokernel of a map of free modules gives natural isomorphisms of chain complexes
\[
 E(e_mU)\cong K_m\otimes_kU,
 \qquad E(Q)\cong\bigoplus_{m\geq1}K_m\otimes_kQ_m.
 \tag{C-tensor-coefficients}\label{Ctensorcoefficients}
\]
We have used exactness and coproduct preservation to prove this tensor description, not to assume commutation of an arbitrary right adjoint with coproducts.

For every pure-injective module $I$, the extension of the counit of $F\dashv G$ across the pure interval map, as in Claim~\ref{C14}, makes $GD(e_mI)$ a retract of its interval tensor. Condition (d) and the resolving property of $F^{(0)}$ give
\[
 K_m\otimes_k I\simeq E(e_mI)\simeq0.
\]
We claim this forces $K_m$ to be pure acyclic. Write $V^+=\operatorname{Hom}_{\mathbb Z}(V,\mathbb Q/\mathbb Z)$. For any pure-injective $k$-module $P$, the character dual $P^+$ is pure-injective, so the preceding vanishing and exact character duality imply
\[
 \operatorname{Hom}_k(K_m,P^{++})
 \cong (K_m\otimes_kP^+)^+
 \quad\text{is acyclic}.
\]
The canonical pure monomorphism $P\to P^{++}$ splits because $P$ is pure-injective. Thus $\operatorname{Hom}_k(K_m,P)$ is acyclic for every pure-injective $P$. As proved explicitly in Claim~\ref{C16}, testing an ordinary injective cogenerator first gives exactness of $K_m$, and testing character duals proves purity of all its cycle inclusions. Hence $K_m$ is pure acyclic. Its tensor product with every module is therefore acyclic. Formula~\eqref{Ctensorcoefficients}, together with exactness of direct sums of modules, proves $E(Q)$ acyclic for every sequence $Q$. For projective disk coefficients this is $L\varphi(GDQ)=0$, which is Claim~\ref{C10}'s remaining gate.

Finally, if $G$ preserves coproducts, then so does $F^{(0)}GD$. In the stated compactness case the adjunction formula
\[
 G(K)(a)=\operatorname{Hom}_{\Ch_{\geq0}(k)}(F(h_a),K)
\]
shows that $G$ preserves coproducts pointwise. Claim~\ref{C1}'s projective disk decomposition proves the asserted characterization of a finitely generated projective complex. These implications establish the sufficient hypotheses without identifying them with the unrestricted corollary assumptions.
\end{proof}

